\subsection{Covering the short unipotent orbit}\label{sec:count-orbit}

In a deep part of the quotient, the lattice filtration has a large gap. The unipotent radical attached to this gap has a small set of representatives modulo its rational points, which allows cuspidality to control the value of a form.

\begin{proposition}[Covering in the chosen short flag]\label{prop:short-orbit-cover}
In Definition~\ref{hyp:cuspidal-dimension}, take \(q\ge1\) and suppose the smallest size \(\tau(x)\) in the canonical lattice filtration is less than \(1/2\). From the filtration in
Lemma~\ref{lem:canonical-short-flag} choose a step with consecutive
size ratio \(u\le(2\tau(x))^{1/q}\). Modulo its rational lattice
stabilizer and the finite level, the short unipotent orbit admits
representatives \(\exp A\exp Y\) in the real coordinates at \(x\),
where
\[
    \|A\|+\|Y\|\le \rho,\qquad
    \rho=\tfrac14 r^2\sqrt f\,g d_E(2\tau(x))^{1/q}.
\]
\end{proposition}

\begin{proof}
\proofstep{Choose a step by the product of ratios.}

Let \(\tau=t_1<t_2<\cdots<t_a<1/2\) be the short sizes
from Lemma~\ref{lem:canonical-short-flag} and put \(t_{a+1}=1/2\).
There are \(a\le q\) ratios since the sum of the short ranks
is at most \(q\). Then
\[
 \prod_{i=1}^a t_i/t_{i+1}=2\tau,\qquad
 \min_i t_i/t_{i+1}\le(2\tau)^{1/a}\le(2\tau)^{1/q}.
\]
Use a pair \(t,t'\) attaining one such ratio \(u\).
Let \(S\) be the saturated isotropic step ending at \(t\),
of rank \(j\). For the middle block \(Z=S^\perp/S\), the extreme sizes are
\[
 t_{\max}(S)=t,\qquad t_{\min}(Z)\ge t',\qquad t_{\min}(S^\vee)=1/t ,
\]
where the inequality for the middle block is needed only when it is nonzero.

\proofstep{Integral Witt coordinates and congruence lattices.}

Use an integral Witt splitting
\(\Lambda=S\oplus Z\oplus S^\vee\).
To obtain it, extend functionals on the saturated summand
\(S\) to \(\Lambda\) and use self-duality to identify
\(\Lambda/S^\perp\) with \(S^\vee\).
Choose a section by projectivity. Its Hermitian form can
be killed by subtracting an \(S\)-valued correction,
since the pairing between \(S\) and this section is perfect.
In detail, choose an integral scalar \(\theta\) with
\(\theta+\bar\theta=1\). It exists since the trace of the
unramified quadratic integral extension is surjective at
each prime, and so globally. Use a correction whose pairing
with the section is \(\theta\) times that section's form.
Indeed, if \(H_0\) is the Hermitian form of the section, the two
cross terms sum to
\[
                  \theta H_0+\bar\theta H_0=H_0 .
\]
Take the orthogonal complement of the resulting hyperbolic
pair to split off \(Z\).

The radical of the stabilizer of \(S\) has a two-step
Lie algebra with coordinates \(X:Z\longrightarrow S\)
and \(Y:S^\vee\longrightarrow S\), where \(Y\) is
skew-Hermitian. The other first-layer block is
\(-X^\dagger\). At the real places use the orthogonal Witt
coordinates. Two such splittings inducing the same maps on
the graded spaces differ by an element of the radical, as
their change of splitting is unitary, block triangular and
the identity on the graded spaces. Conjugation by it acts
trivially on the first-layer quotient and on the center.
Thus the congruence lattice supplies \(2g\operatorname{Hom}(Z,S)\)
in the first block.

For the first layer in the integral splitting, write
\[
 A_X=\begin{pmatrix}0&X&0\\0&0&-X^\dagger\\0&0&0\end{pmatrix},
 \qquad
 \exp(A_X)=
 \begin{pmatrix}1&X&-XX^\dagger/2\\
                 0&1&-X^\dagger\\0&0&1\end{pmatrix}.
\]
For \(X\in2g\operatorname{Hom}(Z,S)\) every off-diagonal
entry is integral and divisible by \(g\). In particular, writing
\(X=2gX_0\) makes the central entry
\[
 -XX^\dagger/2=-2g^2X_0X_0^\dagger
                   \in g\operatorname{Hom}(S^\vee,S).
\]
The center supplies \(g\) times the skew-Hermitian sublattice
of \(\operatorname{Hom}(S^\vee,S)\).

\proofstep{Dual minima of the two coordinate lattices.}

Here are the covering bounds for these lattices. In the
product Hilbert--Schmidt norm the Euclidean dual of
\(\operatorname{Hom}(Z,S)\) identifies with
\(2\mathfrak D_E^{-1}\operatorname{Hom}(S,Z)\), where
\(\mathfrak D_E\) is the absolute different. This follows
from the matrix trace pairing on projective modules and
the factor 2 relating field trace and the real product
pairing over the complex places. Multiplication of a
lattice by the codifferent multiplies every filtration
size by \(d_E^{-1}\), from its ideal norm and the same
exterior-volume normalization. The determinant comparison
for nonzero maps therefore bounds the geometric mean of
the local norms of a nonzero dual matrix below by
\(2d_E^{-1}t'/t\). For a nonzero vector \(B\) in this dual, the quadratic and
geometric means give
\[
 \Big(\sum_v\|B_v\|_{\rm HS}^2\Big)^{1/2}
 \ge\sqrt f\Big(\prod_v\|B_v\|_{\rm HS}\Big)^{1/f}
 \ge 2\sqrt f\,d_E^{-1}t'/t .
\] The first-layer lattice has
real dimension \(2fj(r-2j)\).

For the center, the skew-Hermitian sublattice is primitive
in the full matrix lattice. An integer functional on it
extends to the full lattice, so its Euclidean dual is the
orthogonal skew projection of the full dual. The projection \(B\mapsto (B-B^\dagger)/2\) places the dual of the
central lattice in \(\mathfrak D_E^{-1}\operatorname{Hom}(S,S^\vee)\).
The extreme sizes of \(S\) and its dual thus give
\[
                  \|B\|\ge \sqrt f\,d_E^{-1}/t^2
                         \qquad(B\ne0\text{ in the dual}).
\] The center has real dimension
\(fj^2\).

\proofstep{Cover each layer by transference.}

The dual minimum bounds give a covering radius in each additive layer. Combining these coverings must respect the order of the unipotent extension.


For a full Euclidean lattice of dimension \(d\), its
covering radius times the length of the shortest nonzero
dual vector is at most \(d/2\)
\cite[Theorem~2.2]{banaszczyk}. Apply this to the two
computed dual minima. Including the level \(2g\) in the
first layer and both its adjoint blocks, its radius is at most
\[
 \sqrt2\,(2g)\frac{j(r-2j)\sqrt f\,d_E}{2}\frac{t}{t'}
        \le\frac{r^2\sqrt f\,g d_E}{4\sqrt2}\frac{t}{t'} .
\]
In the center, including level \(g\), the radius is at most
\[
        \frac{g j^2\sqrt f\,d_E}{2}t^2
        \le\frac{r^2\sqrt f\,g d_E}{32}\frac{t}{t'} ,
\]
using \(t t'\le1/4\). First reduce the first block by an
integral congruence lift, then reduce the resulting center,
which already includes the commutator correction.
The two displayed radii can be added after this central correction.
For the gap \(t/t'\le(2\tau)^{1/q}\) this gives
\[
 \|A\|+\|Y\|\le r^2\sqrt f\,g d_E\frac{t}{t'}
       \left(\frac{1}{4\sqrt2}+\frac1{32}\right)
 \le \frac{r^2\sqrt f\,g d_E}{4}(2\tau)^{1/q}=:\rho .
\]

The radical is two-step with vector-group center and
vector-group quotient. Additive approximation at the
finite adeles in these two groups shows that each class
on its adelic orbit modulo the finite level has a real
representative. Apply the real covering just proved to
that representative. A rational left correction from
the lattice stabilizer at level \(g\) gives only a right
level translate at the point, so the value of each projected
function is unchanged. After conjugating into the real
coordinates at \(x\), the representative acts by a product
\(\exp A\exp Y\) with \(\|A\|+\|Y\|\le\rho\).
\end{proof}
